% ============================================================
% Sparse Autoencoders Reveal and Control Gender Bias in
% Vision-Language Model Captioning
% ============================================================
% LaTeX Paper Draft — 3 March 2026
% Target: ACL / EMNLP / NeurIPS style (single-column draft)
% ============================================================

\documentclass[11pt,a4paper]{article}

% ---- Packages ----
\usepackage[utf8]{inputenc}
\usepackage[T1]{fontenc}
\usepackage{amsmath,amssymb,amsfonts}
\usepackage{graphicx}
\usepackage{booktabs}
\usepackage[margin=2.5cm]{geometry}
\usepackage{url}

% Convenience
\newcommand{\modelname}[1]{\textsc{#1}}
\newcommand{\pali}{\modelname{PaLiGemma-3B}}
\newcommand{\qwen}{\modelname{Qwen2-VL-7B}}
\newcommand{\llama}{\modelname{Llama-3.2-11B}}
\newcommand{\llava}{\modelname{LLaVA-1.5-7B}}
\newcommand{\sae}{SAE}
\newcommand{\saes}{SAEs}

\title{Sparse Autoencoders Reveal and Control Gender Bias\\in Vision-Language Model Captioning:\\A Cross-Lingual Mechanistic Interpretability Study}

\author{
  Nour Mubarak \\
  Durham University \\
  \texttt{nour.mubarak@durham.ac.uk}
}

\date{March 2026 — Draft}

\begin{document}
\maketitle

% ============================================================
\begin{abstract}
% ============================================================
Vision-language models (VLMs) generate image captions that often reflect societal gender biases, but the internal mechanisms producing this bias remain poorly understood.
We apply sparse autoencoders (\saes{}) to the internal activations of three VLM architectures---\pali{}, \qwen{}, and \llama{}---to identify and causally intervene on features associated with gender bias in captioning.

Our main findings are threefold.
\textbf{(1)}~Targeted ablation of the top 100 \sae{}-identified gender features (0.6\% of all features) produces a statistically significant change in gender-marked language across all three models, with effect ratios of 1.8--7.1$\times$ relative to random feature ablation controls (25~runs each, 500~images, bootstrap CIs excluding zero).
\textbf{(2)}~The \emph{direction} of the intervention effect depends on model architecture: \pali{} shows a $-16.1\%$ \emph{decrease} in gender terms, while \qwen{} ($+3.95\%$) and \llama{} ($+5.02\%$) show \emph{increases}---suggesting that smaller captioning models encode gender as \emph{production} features, whereas larger instruction-tuned models encode gender through \emph{regulatory/suppressor} features.
\textbf{(3)}~Cross-lingual analysis across English and Arabic reveals near-zero feature overlap (CLBAS~$< 0.015$), demonstrating that gender bias is encoded through language-specific features, with implications for multilingual debiasing.

To our knowledge, this is the first application of \saes{} for mechanistic interpretability of gender bias in VLMs.
Code and data are available at \url{https://github.com/nour-mubarak/mechanistic_intrep}.
\end{abstract}

% ============================================================
\section{Introduction}
\label{sec:intro}
% ============================================================

Vision-language models (VLMs) have demonstrated remarkable capabilities in image captioning, visual question answering, and multimodal reasoning \citep{liu2023visual,bai2023qwen,team2024paligemma}.
However, these models inherit and amplify societal biases present in their training data, particularly gender bias in describing people \citep{hendricks2018women,hirota2022quantifying}.
A man on a skateboard is described as ``a \emph{man} performing tricks,'' while a woman in a similar pose is described as ``a \emph{person} standing on a board''---the model systematically deploys gendered language for some groups while using neutral language for others.

Prior work has \emph{measured} such biases through caption-level statistics \citep{zhao2017men,hendricks2018women} and proposed output-level mitigations such as data augmentation or constrained decoding \citep{zhao2017men}.
However, these approaches treat the model as a black box: they quantify \emph{what} the model does but not \emph{why} or \emph{how} it encodes gender information internally.

Mechanistic interpretability---the study of how neural network components implement specific computations---offers a path toward understanding \emph{how} bias is encoded in model representations.
Sparse autoencoders (\saes{}), recently applied to text-only language models \citep{cunningham2023sparse,bricken2023monosemanticity,templeton2024scaling}, decompose dense neural activations into interpretable, monosemantic features.
If gender bias is encoded through specific features, these features can be identified, analyzed, and surgically ablated to test causal claims.

In this work, we extend \sae{}-based mechanistic interpretability to VLMs for the first time, focusing on gender bias in image captioning across four architectures and two languages.
Our contributions are:

\begin{enumerate}
    \item \textbf{Methodological:} First application of \saes{} for mechanistic interpretability of gender bias in VLMs.
    We train \saes{} on language-model decoder activations from four VLMs (\pali{}, \qwen{}, \llava{}, \llama{}) and identify gender-associated features via differential activation analysis.

    \item \textbf{Empirical:} Causal intervention evidence across three architectures showing that targeted ablation of 100 \sae{} features (0.6\% of total) produces statistically significant changes in gender-marked language---with effect ratios of 1.8--7.1$\times$ relative to random ablation controls.

    \item \textbf{Cross-lingual:} First cross-lingual bias alignment analysis (CLBAS metric) across four VLMs, showing near-zero feature overlap between English and Arabic gender representations---implying that debiasing in one language does not transfer to another.

    \item \textbf{Mechanistic insight:} Discovery that the \emph{direction} of gender-feature influence varies across architectures: smaller captioning models use excitatory gender features (ablation reduces gender terms), while larger instruction-tuned models use inhibitory/regulatory features (ablation increases gender terms).
\end{enumerate}

% ============================================================
\section{Related Work}
\label{sec:related}
% ============================================================

\paragraph{Gender bias in VLMs.}
\citet{hendricks2018women} documented systematic gender bias in image captioning, showing that models describe men with activity-related terms and women with appearance-related terms.
\citet{zhao2017men} proposed data augmentation to reduce bias in visual semantic role labeling.
\citet{hirota2022quantifying} developed metrics for quantifying bias in VLM-generated captions.
These works measure and mitigate bias at the output level but do not explain the internal mechanisms.

\paragraph{Mechanistic interpretability with SAEs.}
\citet{cunningham2023sparse} demonstrated that \saes{} can decompose GPT-2 activations into interpretable, monosemantic features.
\citet{bricken2023monosemanticity} scaled this to larger models, showing that individual features correspond to semantically meaningful concepts.
\citet{templeton2024scaling} trained \saes{} with 34M features on Claude Sonnet, identifying features for safety-relevant concepts.
All prior \sae{} work has focused on text-only models; we extend this to multimodal VLMs.

\paragraph{Causal interpretability.}
Causal tracing \citep{meng2022locating} identifies which model components store factual associations by corrupting and restoring activations.
Activation patching \citep{geiger2021causal} tests causal hypotheses about intermediate representations.
Our approach is complementary: \saes{} decompose activations into interpretable features, and ablation tests whether specific features causally influence output.

% ============================================================
\section{Models and Data}
\label{sec:models}
% ============================================================

\subsection{Models}

We study four VLMs spanning 2.9B--10.6B parameters (Table~\ref{tab:models}).
All models use a vision encoder connected to a language model decoder; we extract activations from the \emph{language model decoder layers} (not the vision encoder), as these are the layers that generate text tokens.

\begin{table}[h]
\centering
\small
\caption{Vision-language models used in this study. $d_{\text{model}}$ is the decoder hidden dimension; $n_{\text{features}}$ is the \sae{} dictionary size (8$\times$ expansion).}
\label{tab:models}
\begin{tabular}{@{}lccccl@{}}
\toprule
\textbf{Model} & \textbf{Params} & $d_{\text{model}}$ & \textbf{Layers} & $n_{\text{features}}$ & \textbf{Intervention} \\
\midrule
\pali{}    & 2.9B  & 2,048 & 18 & 16,384 & L9, L17 \\
\qwen{}    & 7.6B  & 3,584 & 28 & 28,672 & L12 \\
\llava{}   & 7.1B  & 4,096 & 32 & 32,768 & --- \\
\llama{}   & 10.6B & 4,096 & 40 & 32,768 & L20 \\
\bottomrule
\end{tabular}
\end{table}

\paragraph{Layer selection.}
Intervention layers are selected at approximately the 50th percentile of each model's depth, following evidence that middle transformer layers encode the richest semantic features \citep{cunningham2023sparse,templeton2024scaling}.
For \pali{}, L17 (penultimate layer) serves as a negative control.

\subsection{Dataset}

We use a COCO-derived bilingual image captioning dataset containing 8,092 images with 40,455 English--Arabic caption pairs.
Gender labels are derived from caption content using keyword matching (male: 5,562 samples; female: 2,047 samples, reflecting COCO's inherent male bias).

\paragraph{Gender term dictionary.}
We track 16 binary gender terms: \texttt{he, she, him, her, his, hers, man, woman, boy, girl, men, women, boys, girls, male, female}.
Following reviewer feedback, we separately track 18 gender-neutral terms (\texttt{person, people, individual}, etc.) to detect whether ablation shifts output toward neutral descriptions.

% ============================================================
\section{Method}
\label{sec:method}
% ============================================================

Our pipeline has four stages: (1)~activation extraction, (2)~\sae{} training, (3)~gender feature identification, and (4)~causal intervention.

\subsection{Activation Extraction}

For each model, we extract activations from selected decoder layers using PyTorch forward hooks.
Activations are \textbf{mean-pooled} across sequence positions to yield one $\mathbb{R}^{d_{\text{model}}}$ vector per image--caption pair.

\begin{table}[h]
\centering
\small
\caption{Activation extraction details per model.}
\label{tab:extraction}
\begin{tabular}{@{}lccc@{}}
\toprule
\textbf{Model} & \textbf{N samples} & \textbf{Layers extracted} & \textbf{Source} \\
\midrule
\pali{}  & 10,000 & 0,3,6,9,12,15,17 & samples.csv \\
\qwen{}  & 7,609 (EN) / 6,413 (AR) & 0,4,8,12,16,20,24,27 & gender-labeled \\
\llava{} & 5,249 (EN) / 4,449 (AR) & 0,4,8,...,31 & gender-labeled \\
\llama{} & 5,249 (EN) / $\sim$5,000 (AR) & 0,5,10,...,39 & gender-labeled \\
\bottomrule
\end{tabular}
\end{table}

\subsection{SAE Architecture and Training}

We use a vanilla sparse autoencoder with ReLU activation and L1 sparsity:
\begin{align}
    \mathbf{h} &= \text{ReLU}(\mathbf{W}_{\text{enc}} \mathbf{x} + \mathbf{b}_{\text{enc}}) \label{eq:enc} \\
    \hat{\mathbf{x}} &= \mathbf{W}_{\text{dec}} \mathbf{h} + \mathbf{b}_{\text{dec}} \label{eq:dec} \\
    \mathcal{L} &= \text{MSE}(\mathbf{x}, \hat{\mathbf{x}}) + \lambda \|\mathbf{h}\|_1 \label{eq:loss}
\end{align}
where $\mathbf{W}_{\text{enc}} \in \mathbb{R}^{n_{\text{features}} \times d_{\text{model}}}$, $\mathbf{W}_{\text{dec}} \in \mathbb{R}^{d_{\text{model}} \times n_{\text{features}}}$, and $\lambda$ controls sparsity (1$\times$10$^{-4}$ for \pali{}, 5$\times$10$^{-4}$ for \qwen{}).

Training uses AdamW with learning rate $10^{-4}$, batch size 256, 50 epochs with early stopping (patience=10), and a 90/10 train/validation split.
The expansion factor is 8$\times$ throughout.

\subsection{Gender Feature Identification}
\label{sec:feature_id}

Gender-associated features are identified by \textbf{differential activation analysis}:
\begin{enumerate}
    \item Encode all gender-labeled samples through the trained \sae{} encoder
    \item Compute mean feature activation $\bar{a}_j^{\text{male}}$ and $\bar{a}_j^{\text{female}}$ for each feature $j$
    \item Rank features by $|\bar{a}_j^{\text{male}} - \bar{a}_j^{\text{female}}|$
    \item Select top-$k$ features ($k=100$, i.e., 0.6\% of features for \pali{})
\end{enumerate}

\subsection{Causal Intervention}
\label{sec:intervention}

The intervention experiment has three phases:

\begin{enumerate}
    \item \textbf{Baseline:} Generate captions for 500 images (seed=42) with no modification
    \item \textbf{Targeted ablation:} Register a forward hook that zeros the contribution of the $k=100$ gender features during generation
    \item \textbf{Random control:} Repeat 25 times with randomly selected features (same $k$)
\end{enumerate}

\paragraph{Ablation hook.}
We implement two hook variants:

\textbf{Full reconstruction} (used for \pali{}): Replace activations with full \sae{} reconstruction after zeroing targeted features:
\begin{equation}
    \mathbf{x}_{\text{mod}} = \mathbf{W}_{\text{dec}} \cdot (\mathbf{h} \odot \mathbf{m}) + \mathbf{b}_{\text{dec}}
\end{equation}
where $\mathbf{m}$ is a binary mask with $m_j = 0$ for ablated features.

\textbf{Residual ablation} (used for \qwen{}, \llama{}): Subtract only the contribution of targeted features from the original activations:
\begin{equation}
    \mathbf{x}_{\text{mod}} = \mathbf{x} - \sum_{j \in \mathcal{A}} h_j \cdot \mathbf{W}_{\text{dec}}^{(:,j)}
    \label{eq:residual}
\end{equation}
where $\mathcal{A}$ is the set of ablated feature indices.

\paragraph{Justification for two hook variants.}
The full-reconstruction hook works when the \sae{} produces high-fidelity per-token reconstructions.
For \pali{} (cos\_sim = 0.9999), this was the case.
For \qwen{}, the \sae{} was trained on mean-pooled activations but applied to per-token activations during generation; the distribution mismatch caused the full-reconstruction hook to destroy output coherence (producing 0 gender terms in all conditions).
Residual ablation (Eq.~\ref{eq:residual}) avoids this by preserving the original activations and only removing the targeted features' contribution.
Importantly, within each model, the same hook variant is used for both targeted and random conditions, so the targeted-vs.-random comparison remains valid.

\paragraph{Generation.}
All captions are generated with deterministic greedy decoding (\texttt{do\_sample=False}, \texttt{num\_beams=1}, \texttt{max\_new\_tokens=64}) to eliminate sampling variance.

% ============================================================
\section{Statistical Analysis}
\label{sec:stats}
% ============================================================

\paragraph{Primary metric.}
Total gender term count across 500 captions, reported as percentage change from baseline.

\paragraph{Length normalization.}
Gender rate = total gender terms / total tokens, controlling for caption length changes.

\paragraph{Paired statistics.}
For each image $i$, we compute $\Delta_i = \text{count}(\text{ablated}_i) - \text{count}(\text{baseline}_i)$.
We report the mean $\Delta$ with bootstrap 95\% CIs (10,000 resamples) and Wilcoxon signed-rank tests.

\paragraph{One-sided Wilcoxon limitation.}
Our Wilcoxon test uses the one-sided alternative $H_1$: targeted $<$ baseline (testing for decrease).
For models where targeted ablation \emph{increases} gender terms (\qwen{}, \llama{}), the one-sided $p$-value approaches 1.0---this does not indicate absence of effect but rather an effect in the opposite direction.
For these models, significance is established via bootstrap CIs on the paired difference (targeted $-$ random): if the 95\% CI excludes zero, the effect is statistically significant regardless of direction.

\paragraph{Effect specificity.}
$\text{Specificity} = |\text{targeted change\%}| - |\text{random change\%}|$.
Ratio = $|\text{targeted}| / |\text{random}|$; values $> 1$ indicate specificity.

% ============================================================
\section{Results}
\label{sec:results}
% ============================================================

\subsection{Main Intervention Results}

Table~\ref{tab:main_results} presents the core findings across all three models with intervention experiments.

\begin{table*}[t]
\centering
\small
\caption{Intervention results across three VLM architectures. All experiments use 500 images, $k=100$ ablated features, and 25 random control runs. $\Delta$CI is the bootstrap 95\% CI on the paired difference (targeted $-$ random) per image. Significance: CI excludes zero for all three models.}
\label{tab:main_results}
\begin{tabular}{@{}lcccccccc@{}}
\toprule
\textbf{Model} & \textbf{Hook} & \textbf{Baseline} & \textbf{Targeted} & \textbf{Change} & \textbf{Random} & \textbf{Spec.} & \textbf{Ratio} & \textbf{$\Delta$CI (targeted$-$random)} \\
\midrule
\pali{} (L9) & Full recon. & 1,522 & 1,277 & $-16.1\%$ & $-8.7\% \pm 3.9\%$ & $-7.3$ pp & $1.8\times$ & $[-0.274, -0.174]$ \\
\qwen{} (L12) & Residual & 1,315 & 1,367 & $+3.95\%$ & $-0.56\% \pm 1.14\%$ & $+4.5$ pp & $7.1\times$ & $[+0.038, +0.200]$ \\
\llama{} (L20) & Residual & 1,355 & 1,423 & $+5.02\%$ & $-0.84\% \pm 0.76\%$ & $+5.9$ pp & $6.0\times$ & $[+0.078, +0.242]$ \\
\bottomrule
\end{tabular}
\end{table*}

All three models show statistically significant targeted effects: the bootstrap 95\% CI on the paired difference (targeted $-$ random) excludes zero in every case.
Effect specificity ranges from 4.5 to 7.3 percentage points, and effect ratios range from 1.8$\times$ to 7.1$\times$.

\paragraph{Direction divergence.}
The most striking finding is the \emph{direction} of the effect.
\pali{} shows a $-16.1\%$ decrease in gender terms under targeted ablation, while \qwen{} ($+3.95\%$) and \llama{} ($+5.02\%$) show increases.
This suggests two distinct mechanisms of gender encoding:
\begin{itemize}
    \item \pali{} (smaller, captioning-focused): Gender features \emph{produce} gendered output; ablation removes it.
    \item \qwen{}/\llama{} (larger, instruction-tuned): Top-differential features \emph{regulate/suppress} gender; ablation disinhibits it.
\end{itemize}
This is analogous to the distinction between excitatory and inhibitory circuits in neuroscience.

\paragraph{Instruction tuning effect.}
The two instruction-tuned models (\qwen{}, \llama{}) show notably tighter random ablation variance (std = 1.14\%, 0.76\%) compared to \pali{} (std = 3.9\%), suggesting that instruction tuning creates more stable internal representations where random feature perturbation has less impact.

\subsection{Layer Specificity (PaLiGemma)}

\begin{table}[h]
\centering
\small
\caption{Layer specificity analysis for \pali{}. L9 (mid-layer) carries the entire effect; L17 (penultimate) shows no significant change.}
\label{tab:layer_specificity}
\begin{tabular}{@{}lcccc@{}}
\toprule
\textbf{Config} & \textbf{Targeted} & \textbf{Random} & \textbf{Wilcoxon $p$} \\
\midrule
L9 (mid-layer) & $-16.1\%$ & $-8.7\% \pm 3.9\%$ & $1.65 \times 10^{-21}$ \\
L17 (late layer) & $-0.8\%$ & $-4.1\% \pm 5.0\%$ & 0.999 (n.s.) \\
L9 + L17 & $-16.1\%$ & $-8.7\% \pm 3.9\%$ & $1.65 \times 10^{-21}$ \\
\bottomrule
\end{tabular}
\end{table}

Gender features are concentrated in the mid-layer (L9): targeted ablation at L17 produces no significant effect ($p = 0.999$), and combining L9+L17 produces identical results to L9 alone.
This confirms that gender bias is encoded in specific intermediate layers, not distributed uniformly.

\subsection{Per-Term Analysis}

Table~\ref{tab:per_term} shows the most affected terms across models.

\begin{table}[h]
\centering
\small
\caption{Top per-term changes under targeted ablation. \pali{} reduces pronouns; \qwen{}/\llama{} increase explicit nouns and possessives.}
\label{tab:per_term}
\begin{tabular}{@{}lrrr@{}}
\toprule
\textbf{Term} & \textbf{\pali{}} & \textbf{\qwen{}} & \textbf{\llama{}} \\
\midrule
her & $-73.2\%$ & $+4.1\%$ & $+11.9\%$ \\
he & $-22.6\%$ & $+0.2\%$ & $-0.2\%$ \\
his & $-0.4\%$ & $+6.3\%$ & $+66.3\%$ \\
man & $+10.4\%$ & $+53.6\%$ & $+4.8\%$ \\
woman & $+9.3\%$ & $+43.6\%$ & $-4.3\%$ \\
him & $-83.3\%$ & $+90.9\%$ & $+5.3\%$ \\
men & $-33.3\%$ & $+7.5\%$ & $-43.9\%$ \\
son & $+2.8\%$ & $-30.0\%$ & $-40.6\%$ \\
\bottomrule
\end{tabular}
\end{table}

A consistent cross-model pattern emerges: ablation reshuffles \emph{which} gendered terms appear (e.g., pronouns vs.\ nouns) rather than uniformly increasing or decreasing all terms.
For \pali{}, pronouns (``her'', ``him'', ``she'') are dramatically reduced while explicit nouns (``man'', ``woman'') slightly increase---suggesting the ablated features encode pronoun usage patterns rather than pure gender semantics.

\subsection{SAE Quality Metrics}

\begin{table}[h]
\centering
\small
\caption{SAE quality metrics. \pali{} has the highest quality; \llama{} (original) has the lowest. Intervention results are significant across this quality range.}
\label{tab:sae_quality}
\begin{tabular}{@{}lcccc@{}}
\toprule
\textbf{Model} & \textbf{Cos Sim} & \textbf{Expl. Var.} & \textbf{L0} & \textbf{Dead \%} \\
\midrule
\pali{}  & 0.9999 & 99.8\% & 7,992 & 51.2\% \\
\qwen{}  & 0.9965 & 66.4\% & 2,049 & 71.6\% \\
\llama{} & 0.9956 & 36.6\% & 344 & 98.6\% \\
\bottomrule
\end{tabular}
\end{table}

SAE quality varies substantially across models.
Notably, \llama{} achieves significant intervention results (6.0$\times$ ratio, CI excludes zero) despite having the lowest-quality \sae{} (36.6\% explained variance, 98.6\% dead features).
This suggests the intervention is robust to \sae{} quality and may strengthen with better-trained \saes{}.

\subsection{Cross-Lingual Analysis}

\begin{table}[h]
\centering
\small
\caption{Cross-lingual bias alignment. CLBAS $\approx 0$ indicates gender is encoded through language-specific features, not shared cross-lingual representations.}
\label{tab:clbas}
\begin{tabular}{@{}lccc@{}}
\toprule
\textbf{Model} & \textbf{CLBAS} & \textbf{EN Probe} & \textbf{AR Probe} \\
\midrule
\pali{}  & 0.0004 & 0.810 & 0.812 \\
\qwen{}  & 0.004 & 0.918 & 0.903 \\
\llava{} & 0.015 & 0.963 & 0.899 \\
\llama{} & 0.004 & 0.994 & 0.985 \\
\bottomrule
\end{tabular}
\end{table}

All models encode gender in linearly decodable features (probe accuracy $> 0.80$), but cross-lingual feature overlap is near-zero (CLBAS $< 0.015$, Jaccard overlap $< 3\%$ at all layers).
This means English and Arabic gender biases are encoded through \textbf{entirely separate feature sets}---debiasing in English would not automatically reduce bias in Arabic.

% ============================================================
\section{Discussion}
\label{sec:discussion}
% ============================================================

\paragraph{Excitatory vs.\ inhibitory gender features.}
The most important finding is the direction divergence across architectures.
\pali{} (a simpler, captioning-focused model) appears to encode gender through \emph{excitatory} features: these features directly contribute gendered language to the output, so ablation reduces it.
\qwen{} and \llama{} (larger, instruction-tuned) appear to use \emph{inhibitory/regulatory} features: these features modulate or suppress gendered output, so ablation releases it.

This finding has practical implications: a na\"ive debiasing approach that ablates top-differential features would \emph{reduce} gender bias in \pali{} but \emph{increase} it in \qwen{}/\llama{}.
Architecture-specific analysis is essential before deploying \sae{}-based debiasing interventions.

\paragraph{Hook method confound.}
We acknowledge that the direction difference could partly reflect the different hook methods (full reconstruction vs.\ residual ablation).
However, both \qwen{} and \llama{} use the same residual hook and show the same direction, while the full-reconstruction hook \emph{destroyed} \qwen{}'s output entirely---suggesting the direction difference is primarily architectural.
Running \pali{} with the residual hook would disambiguate this (see Limitations).

\paragraph{Cross-lingual implications.}
The near-zero CLBAS scores across all models challenge the assumption that multilingual models share bias representations across languages.
Even \qwen{} (natively multilingual) shows CLBAS = 0.004.
This implies that multilingual fairness requires language-specific debiasing interventions.

\paragraph{Mean-pooling trade-off.}
Our \saes{} are trained on mean-pooled activations but applied to per-token activations during generation.
For \pali{}, this works because the \sae{} has near-perfect reconstruction on pooled data.
For \qwen{}/\llama{}, the mismatch necessitated residual ablation (\Cref{eq:residual}).
Training per-token \saes{} would avoid this mismatch but at significantly higher computational cost.

% ============================================================
\section{Limitations}
\label{sec:limitations}
% ============================================================

\begin{enumerate}
    \item \textbf{SAE training scale:} \saes{} trained on 5K--10K samples per model; larger training sets may yield better features.
    \item \textbf{Binary gender:} We track binary gender terms only; non-binary terms are tracked separately but not used as the primary metric.
    \item \textbf{Hook inconsistency:} \pali{} uses full reconstruction; \qwen{}/\llama{} use residual ablation. Running \pali{} with residual ablation would test whether the direction difference is architectural or methodological.
    \item \textbf{SAE quality variance:} \llama{}'s \sae{} has only 36.6\% explained variance. The tuned \sae{} (99.9\% EV) exists but was not used for intervention.
    \item \textbf{No human evaluation:} We measure gender term counts but do not evaluate caption fluency or semantic accuracy after ablation.
    \item \textbf{Feature polysemanticity:} We cannot claim the ablated features \emph{exclusively} encode gender; they may be polysemantic.
\end{enumerate}

% ============================================================
\section{Conclusion}
\label{sec:conclusion}
% ============================================================

We have demonstrated that sparse autoencoders can identify and causally intervene on gender bias features in VLM internal representations.
Across three architectures, targeted ablation of just 0.6\% of \sae{} features produces statistically significant changes in gender-marked captioning language, with effect ratios of 1.8--7.1$\times$ relative to random controls.
The discovery that different architectures use fundamentally different gender encoding mechanisms---excitatory vs.\ inhibitory---provides new mechanistic insight into how VLMs process social attributes and has practical implications for bias mitigation strategies.

% ============================================================
% References
% ============================================================
\bibliographystyle{plain}

\begin{thebibliography}{99}

\bibitem[Bai et~al.(2023)]{bai2023qwen}
J.~Bai, S.~Bai, S.~Yang, et~al.
\newblock Qwen-VL: A versatile vision-language model for understanding, localization, text reading, and beyond.
\newblock \emph{arXiv preprint arXiv:2308.12966}, 2023.

\bibitem[Bricken et~al.(2023)]{bricken2023monosemanticity}
T.~Bricken, A.~Templeton, J.~Batson, et~al.
\newblock Towards monosemanticity: Decomposing language models with dictionary learning.
\newblock \emph{Anthropic Research}, 2023.

\bibitem[Cunningham et~al.(2023)]{cunningham2023sparse}
H.~Cunningham, A.~Ewart, L.~Riggs, R.~Huben, and L.~Sharkey.
\newblock Sparse autoencoders find highly interpretable features in language models.
\newblock \emph{arXiv preprint arXiv:2309.08600}, 2023.

\bibitem[Geiger et~al.(2021)]{geiger2021causal}
A.~Geiger, H.~Lu, T.~Icard, and C.~Potts.
\newblock Causal abstractions of neural networks.
\newblock In \emph{NeurIPS}, 2021.

\bibitem[Hendricks et~al.(2018)]{hendricks2018women}
L.~A.~Hendricks, K.~Burns, K.~Saenko, T.~Darrell, and A.~Rohrbach.
\newblock Women also snowboard: Overcoming bias in captioning models.
\newblock In \emph{ECCV}, 2018.

\bibitem[Hirota et~al.(2022)]{hirota2022quantifying}
Y.~Hirota, Y.~Nakashima, and N.~Garcia.
\newblock Quantifying societal bias amplification in image captioning.
\newblock In \emph{CVPR}, 2022.

\bibitem[Liu et~al.(2023)]{liu2023visual}
H.~Liu, C.~Li, Q.~Wu, and Y.~J.~Lee.
\newblock Visual instruction tuning.
\newblock In \emph{NeurIPS}, 2023.

\bibitem[Meng et~al.(2022)]{meng2022locating}
K.~Meng, D.~Bau, A.~Andonian, and Y.~Belinkov.
\newblock Locating and editing factual associations in {GPT}.
\newblock In \emph{NeurIPS}, 2022.

\bibitem[Team et~al.(2024)]{team2024paligemma}
PaLiGemma Team.
\newblock PaLiGemma: A versatile 3B VLM for transfer.
\newblock \emph{arXiv preprint arXiv:2407.07726}, 2024.

\bibitem[Templeton et~al.(2024)]{templeton2024scaling}
A.~Templeton, T.~Conerly, J.~Marcus, et~al.
\newblock Scaling monosemanticity: Extracting interpretable features from Claude 3 Sonnet.
\newblock \emph{Anthropic Research}, 2024.

\bibitem[Zhao et~al.(2017)]{zhao2017men}
J.~Zhao, T.~Wang, M.~Yatskar, V.~Ordonez, and K.-W.~Chang.
\newblock Men also like shopping: Reducing gender bias amplification using corpus-level constraints.
\newblock In \emph{EMNLP}, 2017.

\end{thebibliography}

% ============================================================
\appendix
\section{Additional Per-Term Analysis}
\label{app:per_term}

Complete per-term breakdown for all models is available in the supplementary materials and at \url{https://github.com/nour-mubarak/mechanistic_intrep}.

\section{SAE Training Details}
\label{app:sae_training}

\begin{table}[h]
\centering
\small
\caption{SAE training hyperparameters.}
\begin{tabular}{@{}ll@{}}
\toprule
\textbf{Parameter} & \textbf{Value} \\
\midrule
Expansion factor & 8$\times$ \\
L1 coefficient ($\lambda$) & $10^{-4}$ (\pali{}), $5 \times 10^{-4}$ (\qwen{}) \\
Learning rate & $10^{-4}$ \\
Optimizer & AdamW \\
Batch size & 256 \\
Epochs & 50 (early stopping, patience=10) \\
Warmup steps & 1,000 \\
Train/Val split & 90/10 \\
\bottomrule
\end{tabular}
\end{table}

\section{Random Ablation Distributions}
\label{app:random_dist}

See supplementary Figure~S2 for the distribution of all 25 random ablation runs per model, showing that random effects are centered near zero while targeted effects fall well outside the distribution.

\end{document}
